%%%%%%%%%%%%%%%%
%%% exod 34 %%%
%%%%%%%%%%%%%%%%

\Chapter{34}

\Verse{1}
{
	\hJ{וַיֹּאמֶר יְהֹוָה אֶל מֹשֶׁה פְּסל לְךָ שְׁנֵי לֻחֹת אֲבָנִים}
	\hRJE{כָּרִאשֹׁנִים וְכָתַבְתִּי עַל הַלֻּחֹת אֶת הַדְּבָרִים אֲשֶׁר הָיוּ עַל הַלֻּחֹת הָרִאשֹׁנִים אֲשֶׁר שִׁבַּרְתָּ׃}
}
{
	\eJ{And YHWH said to Moses, “Carve two tablets of stones}
	\eRJE{like the first ones, and I’ll write on the tablets the
		words that were on the first tablets, which you shattered.}
}
{
	\fC{REF: The tablets in E are shattered. Now comes the J account of the tablets. In the combined JE text, it would be awkward to picture God just commanding Moses to make some tablets, as if there were no history to this matter, so RJE adds the explanation that these are a replacement for the earlier tablets that were shattered. In E the tablets that are shattered are never said to be replaced. This suggests that, according to E, the ark that is housed in the Temple in Judah to the south either contains broken tablets or no tablets at all. As in other places, the northern Israel source E has clashing religious symbols from those of the southern kingdom of Judah. See the Collection of Evidence, pp. 19, 21.}
}

\Verse{2}
{
	\hJ{וֶהְיֵה נָכוֹן לַבֹּקֶר וְעָלִיתָ בַבֹּקֶר אֶל הַר סִינַי וְנִצַּבְתָּ לִי שָׁם עַל רֹאשׁ הָהָר׃}
}
{
	\eJ{
		And be ready for the morning, and you shall go up in the
		morning to Mount Sinai and present yourself to me there on
		the top of the mountain.
	}
}
{
}

\Verse{3}
{
	\hJ{וְאִישׁ לֹא יַעֲלֶה עִמָּךְ וְגַם אִישׁ אַל יֵרָא בְּכל הָהָר גַּם הַצֹּאן וְהַבָּקָר אַל יִרְעוּ אֶל מוּל הָהָר הַהוּא׃}
}
{
	\eJ{
		And no man shall go up with you, and also let no man be
		seen in all of the mountain. Also let the flock and the
		oxen not feed opposite that mountain.”
	}
}
{
}

\Verse{4}
{
	\hJE{וַיִּפְסֹל שְׁנֵי לֻחֹת אֲבָנִים}
	\hRJE{כָּרִאשֹׁנִים}
	\hJE{וַיַּשְׁכֵּם מֹשֶׁה בַבֹּקֶר וַיַּעַל אֶל הַר סִינַי כַּאֲשֶׁר צִוָּה יְהֹוָה אֹתוֹ וַיִּקַּח בְּיָדוֹ שְׁנֵי לֻחֹת אֲבָנִים׃}
}
{
	\eJE{And he carved two tablets of stones}
	\eRJE{like the first ones,}
	\eJE{and Moses got up early in the morning and went up to Mount
		Sinai as YHWH had commanded him, and he took in his hand
		two tablets of stones.
	}
}
{
}

\Verse{5}
{
	\hJ{וַיֵּרֶד יְהֹוָה בֶּעָנָן וַיִּתְיַצֵּב עִמּוֹ שָׁם וַיִּקְרָא בְשֵׁם יְהֹוָה׃}
}
{
	\eJ{
		And YHWH came down in a cloud and stood with him there,
		and he invoked the name YHWH.
	}
}
{
}

\Verse{6}
{
	\hJ{וַיַּעֲבֹר יְהֹוָה עַל פָּנָיו וַיִּקְרָא יְהֹוָה יְהֹוָה אֵל רַחוּם וְחַנּוּן אֶרֶךְ אַפַּיִם וְרַב חֶסֶד וֶאֱמֶת׃}
}
{
	\eJ{
		And YHWH passed in front of him and called, YHWH, YHWH,
		merciful and gracious {\God}, slow to anger and abounding in
		kindness and faithfulness,
	}
}
{
%	\footnote{
%		and called, YHWH, YHWH. This is the traditional understanding: that the
%		name of God is repeated. But the verse may also be read thus: “and YHWH
%		called, ‘YHWH, merciful and gracious God … ’”
%	}
	\footnote{\fC{REF: This famous formula in J emphasizes the merciful over the just side of God: mercy, grace, kindness. As noted in the Collection of Evidence (p. 12), P never uses these words or several other words relating to mercy. P rather emphasizes the just side of God. This is an important example of the pervasive way in which the Bible became more than the sum of its parts when the Redactor combined the sources. J (and E and D) emphasized the merciful side of God; P emphasized the just side. The final version of the united Torah now brings the two sides together in a new balance, conveying a picture of God who is torn between His justice and His mercy—which has been a central element of the conception of God in Judaism and Christianity ever since.}}
}

\Verse{7}
{
	\hJ{נֹצֵר חֶסֶד לָאֲלָפִים נֹשֵׂא עָוֺן וָפֶשַׁע וְחַטָּאָה וְנַקֵּה לֹא יְנַקֶּה פֹּקֵד עֲוֺן אָבוֹת עַל בָּנִים וְעַל בְּנֵי בָנִים עַל שִׁלֵּשִׁים וְעַל רִבֵּעִים׃}
}
{
	\eJ{
		keeping kindness for thousands, bearing crime and offense and sin;
		though not making one innocent: reckoning fathers’ crime on children
		and on children’s children, on third generations and on fourth
		generations.
	}
}
{
%	\footnote{
%		34:6–7. merciful, gracious, slow to anger, kindness, faithfulness,
%		bearing crime and offense and sin. This is possibly the most
%		repeated and quoted formula in the Tanak (Num 14:18–19; Jon 4:2;
%		Joel 2:13; Mic 7:18; Pss 86:15; 103:8; 145:8; 2 Chr 30:9; Neh
%		9:17,31). The Torah never says what the essence of {\God} is, in
%		contrast to the pagan gods. Baal is the storm wind, Dagon is grain,
%		Shamash is the sun. But what is YHWH? This formula, expressed in the
%		moment of the closest revelation any human has of {\God} in the
%		Bible, is the closest the Torah comes to describing the nature of
%		{\God}. Although humans are not to know what the essence is, they
%		can know what are the marks of the divine personality: mercy, grace.
%		In eight (or nine) different ways we are told of {\God}’s
%		compassion. The last line of the formula (“though not making one
%		innocent”) conveys that this does not mean that one can just get
%		away with anything; there is still justice. But the formula clearly
%		places the weight on divine mercy over divine justice, and it never
%		mentions divine anger. Those who speak of the “Old Testament {\God}
%		of wrath” focus disproportionately on the episodes of anger in the
%		Bible and somehow lose this crucial passage and the hundreds of
%		times that the divine mercy functions in the Hebrew Bible.
%
%		Footnote 34:7. on fourth generations. The list covers five
%		generations: from parents to the fourth generation of their
%		offspring (great-great-grandchildren). Why? First, five generations
%		constitutes the likely maximum of living acquaintance. People
%		occasionally live to see their great-great-grandchildren, but
%		great-great-great is practically unheard of. Second, Israelite
%		rock-cut tombs prior to the seventh century were multi chambered,
%		with room for at least four generations of offspring. The ancestral
%		bond was thus thought to go at least that far back. Third,
%		psychologically one can observe traits persevering through that many
%		generations. I have personally observed ongoing dynamics within my
%		family through four generations. This does not mean that an
%		individual’s bad deed will be duplicated by his or her children and
%		grandchildren. But it may recognize that such deeds have
%		consequences, for better or worse (pride or embarrassment, stigmas,
%		reactions, conscious or unconscious imitation), that persist through
%		generations.
%	}
}

\Verse{8}
{
	\hJ{וַיְמַהֵר מֹשֶׁה וַיִּקֹּד אַרְצָה וַיִּשְׁתָּחוּ׃}
}
{
	\eJ{And Moses hurried and knelt to the ground and bowed,}
}
{
}

\Verse{9}
{
	\hJ{וַיֹּאמֶר אִם נָא מָצָאתִי חֵן בְּעֵינֶיךָ אֲדֹנָי יֵלֶךְ נָא אֲדֹנָי בְּקִרְבֵּנוּ כִּי עַם קְשֵׁה עֹרֶף הוּא וְסָלַחְתָּ לַעֲוֺנֵנוּ וּלְחַטָּאתֵנוּ וּנְחַלְתָּנוּ׃}
}
{
	\eJ{
		and he said, “If I’ve found favor in your eyes, my Lord,
		may my Lord go among us, because it is a stiff-necked
		people, and forgive our crime and our sin, and make us
		your legacy.”
	}
}
{
}

\Verse{10}
{
	\hJ{וַיֹּאמֶר הִנֵּה אָנֹכִי כֹּרֵת בְּרִית נֶגֶד כּל עַמְּךָ אֶעֱשֶׂה נִפְלָאֹת אֲשֶׁר לֹא נִבְרְאוּ בְכל הָאָרֶץ וּבְכל הַגּוֹיִם וְרָאָה כל הָעָם אֲשֶׁר אַתָּה בְקִרְבּוֹ אֶת מַעֲשֵׂה יְהֹוָה כִּי נוֹרָא הוּא אֲשֶׁר אֲנִי עֹשֶׂה עִמָּךְ׃}
}
{
	\eJ{
		And He said, “Here, I’m making a covenant. Before all your
		people I’ll do wonders that haven’t been created in all
		the earth and among all the nations; and all the people
		whom you’re among will see YHWH’s deeds, because that
		which I’m doing with you is awesome.
	}
}
{
}

\Verse{11}
{
	\hJ{שְׁמר לְךָ אֵת אֲשֶׁר אָנֹכִי מְצַוְּךָ הַיּוֹם הִנְנִי גֹרֵשׁ מִפָּנֶיךָ אֶת הָאֱמֹרִי וְהַכְּנַעֲנִי וְהַחִתִּי וְהַפְּרִזִּי וְהַחִוִּי וְהַיְבוּסִי׃}
}
{
	\eJ{
		Watch yourself regarding what I command you today. Here,
		I’m driving out from before you the Amorite and the
		Canaanite and the Hittite and the Perizzite and the Hivite
		and the Jebusite.
	}
}
{
}

\Verse{12}
{
	\hJ{הִשָּׁמֶר לְךָ פֶּן תִּכְרֹת בְּרִית לְיוֹשֵׁב הָאָרֶץ אֲשֶׁר אַתָּה בָּא עָלֶיהָ פֶּן יִהְיֶה לְמוֹקֵשׁ בְּקִרְבֶּךָ׃}
}
{
	\eJ{
		Be watchful of yourself that you don’t make a covenant
		with the resident of the land onto which you’re coming,
		that he doesn’t become a trap among you.
	}
}
{
}

\Verse{13}
{
	\hJ{כִּי אֶת מִזְבְּחֹתָם תִּתֹּצוּן וְאֶת מַצֵּבֹתָם תְּשַׁבֵּרוּן וְאֶת אֲשֵׁרָיו תִּכְרֹתוּן׃}
}
{
	\eJ{
		But you shall demolish their altars and shatter their
		pillars and cut down their Asherahs.
	}
}
{
}

\Verse{14}
{
	\hJ{כִּי לֹא תִשְׁתַּחֲוֶה לְאֵל אַחֵר כִּי יְהֹוָה קַנָּא שְׁמוֹ אֵל קַנָּא הוּא׃}
}
{
	\eJ{
		For you shall not bow to another god—because YHWH: His
		name is Jealous, He is a jealous {\God}—
	}
}
{
%	\footnote{
%		a jealous God. See the comment on Exod 20:5.
%	}
}

\Verse{15}
{
	\hJ{פֶּן תִּכְרֹת בְּרִית לְיוֹשֵׁב הָאָרֶץ וְזָנוּ אַחֲרֵי אֱלֹהֵיהֶם וְזָבְחוּ לֵאלֹהֵיהֶם וְקָרָא לְךָ וְאָכַלְתָּ מִזִּבְחוֹ׃}
}
{
	\eJ{
		that you not make a covenant with the resident of the
		land, and they will prostitute themselves after their gods
		and sacrifice to their gods, and he will call to you, and
		you will eat from his sacrifice.
	}
}
{
}

\Verse{16}
{
	\hJ{וְלָקַחְתָּ מִבְּנֹתָיו לְבָנֶיךָ וְזָנוּ בְנֹתָיו אַחֲרֵי אֱלֹהֵיהֶן וְהִזְנוּ אֶת בָּנֶיךָ אַחֲרֵי אֱלֹהֵיהֶן׃}
}
{
	\eJ{
		And you will take some of his daughters for your sons, and
		his daughters will prostitute themselves after their gods
		and cause your sons to prostitute themselves after their
		gods.
	}
}
{
%	\footnote{
%		you will take some of his daughters for your sons. The Septuagint adds,
%		“and you will give some of your daughters to their sons,” which was
%		probably lost in the Masoretic Text because a scribe’s eye jumped from
%		one occurrence of the word “sons” to the next. Marrying persons from
%		other groups happens with both sexes.
%	}
}

\Verse{17}
{
	\hJ{אֱלֹהֵי מַסֵּכָה לֹא תַעֲשֶׂה לָּךְ׃}
}
{
	\eJ{You shall not make molten gods for yourself.}
}
{
%	\footnote{
%		molten gods. Only after the golden calf incident is the commandment
%		added: Don’t make molten gods (masskh)! If there was any doubt about
%		their permissibility before, there is none now. In the wake of that
%		event, no such statue is ever to be made again. This is notable in
%		connection with the point I made above, that the golden calf is a throne
%		dais for deities, because there will be two golden cherubs in the Temple
%		of Solomon. The cherubs are the throne dais of YHWH, but they are not
%		molten; they are rather made of wood and then gold-plated (1 Kings
%		6:23,28). They are therefore not criticized. But when the northern
%		kingdom of Israel establishes two molten golden calves (1 Kings 12), the
%		kings of Israel are condemned for having them (2 Kings 10:29; 17:22–23).
%	}
}

\Verse{18}
{
	\hJ{אֶת חַג הַמַּצּוֹת תִּשְׁמֹר שִׁבְעַת יָמִים תֹּאכַל מַצּוֹת אֲשֶׁר צִוִּיתִךָ לְמוֹעֵד חֹדֶשׁ הָאָבִיב כִּי בְּחֹדֶשׁ הָאָבִיב יָצָאתָ מִמִּצְרָיִם׃}
}
{
	\eJ{
		You shall observe the Festival of Unleavened Bread. Seven
		days you shall eat unleavened bread, which I commanded
		you, at the appointed time, the month of Abib; because in
		the month of Abib you went out of Egypt.
	}
}
{
}

\Verse{19}
{
	\hJ{כּל פֶּטֶר רֶחֶם לִי וְכל מִקְנְךָ תִּזָּכָר פֶּטֶר שׁוֹר וָשֶׂה׃}
}
{
	\eJ{
		Every first birth of a womb is mine, and all your animals
		that have a male first birth, ox or sheep.
	}
}
{
}

\Verse{20}
{
	\hJ{וּפֶטֶר חֲמוֹר תִּפְדֶּה בְשֶׂה וְאִם לֹא תִפְדֶּה וַעֲרַפְתּוֹ כֹּל בְּכוֹר בָּנֶיךָ תִּפְדֶּה וְלֹא יֵרָאוּ פָנַי רֵיקָם׃}
}
{
	\eJ{
		And you shall redeem an ass’s first birth with a sheep,
		and if you do not redeem it then you shall break its neck.
		You shall redeem every firstborn of your sons. And none
		shall appear before me empty-handed.
	}
}
{
}

\Verse{21}
{
	\hJ{שֵׁשֶׁת יָמִים תַּעֲבֹד וּבַיּוֹם הַשְּׁבִיעִי תִּשְׁבֹּת בֶּחָרִישׁ וּבַקָּצִיר תִּשְׁבֹּת׃}
}
{
	\eJ{
		Six days you shall work, and in the seventh day you shall
		cease. In plowing time and in harvest, you shall cease.
	}
}
{
}

\Verse{22}
{
	\hJ{וְחַג שָׁבֻעֹת תַּעֲשֶׂה לְךָ בִּכּוּרֵי קְצִיר חִטִּים וְחַג הָאָסִיף תְּקוּפַת הַשָּׁנָה׃}
}
{
	\eJ{
		And you shall make a Festival of Weeks, of the firstfruits
		of the wheat harvest, and the Festival of Gathering at the
		end of the year.
	}
}
{
}

\Verse{23}
{
	\hJ{שָׁלֹשׁ פְּעָמִים בַּשָּׁנָה יֵרָאֶה כּל זְכוּרְךָ אֶת פְּנֵי הָאָדֹן יְהֹוָה אֱלֹהֵי יִשְׂרָאֵל׃}
}
{
	\eJ{
		Three times in the year every one of your males shall
		appear before the Lord YHWH, {\God} of Israel.
	}
}
{
}

\Verse{24}
{
	\hJ{כִּי אוֹרִישׁ גּוֹיִם מִפָּנֶיךָ וְהִרְחַבְתִּי אֶת גְּבֻלֶךָ וְלֹא יַחְמֹד אִישׁ אֶת אַרְצְךָ בַּעֲלֹתְךָ לֵרָאוֹת אֶת פְּנֵי יְהֹוָה אֱלֹהֶיךָ שָׁלֹשׁ פְּעָמִים בַּשָּׁנָה׃}
}
{
	\eJ{
		For I shall dispossess nations before you and widen your
		border, and no man will covet your land while you are
		going up to appear before YHWH, your {\God}, three times in
		the year.
	}
}
{
}

\Verse{25}
{
	\hJ{לֹא תִשְׁחַט עַל חָמֵץ דַּם זִבְחִי וְלֹא יָלִין לַבֹּקֶר זֶבַח חַג הַפָּסַח׃}
}
{
	\eJ{
		You shall not offer the blood of my sacrifice on leavened
		bread. And the sacrifice of the Festival of Passover shall
		not remain until the morning.
	}
}
{
%	\footnote{
%		You shall not offer the blood of my sacrifice on leavened bread. Leaven
%		is considered inappropriate for sacrifices. See the comment on Exod
%		29:2.
%	}
}

\Verse{26}
{
	\hJ{רֵאשִׁית בִּכּוּרֵי אַדְמָתְךָ תָּבִיא בֵּית יְהֹוָה אֱלֹהֶיךָ לֹא תְבַשֵּׁל גְּדִי בַּחֲלֵב אִמּוֹ׃}
}
{
	\eJ{
		You shall bring the first of the firstfruits of your land
		to the house of YHWH, your {\God}. You shall not cook a kid
		in its mother’s milk.
	}
}
{
}

\Verse{27}
{
	\hJ{וַיֹּאמֶר יְהֹוָה אֶל מֹשֶׁה כְּתב לְךָ אֶת הַדְּבָרִים הָאֵלֶּה כִּי עַל פִּי הַדְּבָרִים הָאֵלֶּה כָּרַתִּי אִתְּךָ בְּרִית וְאֶת יִשְׂרָאֵל׃}
}
{
	\eJ{
		And YHWH said to Moses, “Write these words for yourself,
		because I’ve made a covenant with you and with Israel
		based on these words.”
	}
}
{
	\footnote{\fC{REF: This famous formula in J emphasizes the merciful over the just side of God: mercy, grace, kindness. As noted in the Collection of Evidence (p. 12), P never uses these words or several other words relating to mercy. P rather emphasizes the just side of God. This is an important example of the pervasive way in which the Bible became more than the sum of its parts when the Redactor combined the sources. J (and E and D) emphasized the merciful side of God; P emphasized the just side. The final version of the united Torah now brings the two sides together in a new balance, conveying a picture of God who is torn between His justice and His mercy—which has been a central element of the conception of God in Judaism and Christianity ever since.}}
}

\Verse{28}
{
	\hJ{וַיְהִי שָׁם עִם יְהֹוָה אַרְבָּעִים יוֹם וְאַרְבָּעִים לַיְלָה לֶחֶם לֹא אָכַל וּמַיִם לֹא שָׁתָה וַיִּכְתֹּב עַל הַלֻּחֹת אֵת דִּבְרֵי הַבְּרִית עֲשֶׂרֶת הַדְּבָרִים׃}
}
{
	\eJ{
		And he was there with YHWH forty days and forty nights. He
		did not eat bread, and he did not drink water. And he
		wrote on the tablets the words of the covenant, the Ten
		Commandments.
	}
}
{
%	\footnote{
%		the Ten Commandments. Hebrew ‘eret haddbrîm, literally, “the ten
%		things.” This is what the Ten Commandments are called in Hebrew.
%		Footnote 34:28. the Ten Commandments. The second set of the commandments
%		appears here in vv. 14–26. Three of them are similar to the commandments
%		that appear in Exodus 20: the commandment against bowing to other gods
%		(34:14–16), the commandment against molten gods (v. 17), and the
%		commandment to cease work on the seventh day (v. 21). The other seven
%		are different from the Ten Commandments that God speaks aloud over
%		Sinai. In critical biblical scholarship we understand these two versions
%		of the Decalogue to come from two different ancient sources. But how are
%		we to understand them in the final form of the Torah? The answer may lie
%		in a second contradiction: In the first verse of this chapter God tells
%		Moses that “I’ll write on the tablets the words that were on the first
%		tablets.” But now God tells Moses, “Write these words for yourself”
%		(34:27). Perhaps we should understand this to mean that God writes the
%		words on one side of the tablets, and Moses writes the words of the
%		second set of commandments on the other side. As is commonly noted, the
%		majority of the first set are ethical commandments, involving relations
%		between humans and other humans: don’t murder, don’t steal … The second
%		set are mainly ritual commandments: observe the holidays, redeem the
%		firstborn, don’t sacrifice with leaven … The two sets are thus
%		complementary, involving the two essential kinds of commandments:
%		relations between humans and humans, and relations between humans and
%		God. Another explanation: We may understand this to mean that God writes
%		the words of the Ten Commandments on both the first and the second set
%		of tablets, while Moses writes these other matters on another document.
%	}
}

\Verse{29}
{
	\hP{וַיְהִי בְּרֶדֶת מֹשֶׁה מֵהַר סִינַי וּשְׁנֵי לֻחֹת הָעֵדֻת בְּיַד מֹשֶׁה בְּרִדְתּוֹ מִן הָהָר וּמֹשֶׁה לֹא יָדַע כִּי קָרַן עוֹר פָּנָיו בְּדַבְּרוֹ אִתּוֹ׃}
}
{
	\eP{
		And it was when Moses was coming down from Mount Sinai,
		and the two tablets of the Testimony were in Moses’ hand
		when he was coming down from the mountain. And Moses had
		not known that the skin of his face was transformed when
		He was speaking with him.
	}
}
{
%	\footnote{
%		transformed. Coming out of the fiery top of the mountain, and back from
%		his once-in-human-history encounter with God, Moses is transformed in
%		some way that is unfortunately obscured in a difficult Hebrew passage
%		(34:30,35). It has often been understood to mean that Moses’ face beams
%		light, and it has been erroneously visualized in numerous artistic
%		depictions as a horned Moses, the most famous of which is Michelangelo’s
%		Moses. William Propp has argued persuasively that it more probably means
%		that Moses’ face is in some way disfigured (from the fire? from the
%		experience of encountering the deity?). Whatever it is about Moses’
%		skin, though, it is a marker, from this point on in the narrative, of
%		Moses’ exceptional position. After having beheld God, he has himself
%		become fearful for other humans to behold. For the rest of the narrative
%		in Exodus (and in the next three books of the Hebrew Bible), he is to be
%		pictured wearing a veil. He removes it only when speaking with the deity
%		and when revealing the divine commandments to the people. Moses’
%		continuing communication with God, moreover, is to take place inside the
%		precincts of the Tabernacle, a place that is forbidden to the lay
%		population. That is, the book of Exodus introduces a series of layers in
%		a channel to God: Moses’ veil, then the Tabernacle, which is a series of
%		tents within tents (the prket, the miškn, the ’hel, the leather cover;
%		see the comment on 26:33) into which Moses enters and from which he
%		returns with the word of God. Interestingly, the layers do not appear to
%		function as barriers between the people and their God in this narrative,
%		but, quite the contrary, there is more of a feeling of connection to the
%		deity via this channel. The veil that hides Moses’ face is, at the same
%		time, a visible indicator of the proximity of deity. The Tabernacle is
%		the visible shrine, confirmed by the appearance of YHWH’s cloud and
%		glory, that marks God’s presence among the people. The fabrics of the
%		veil and the tents give a tangibility to the channel of communication
%		with the divine, and at the same time they preserve the quality of
%		mystery that we have seen to surround the knowledge of YHWH.
%	}
}

\Verse{30}
{
	\hP{וַיַּרְא אַהֲרֹן וְכל בְּנֵי יִשְׂרָאֵל אֶת מֹשֶׁה וְהִנֵּה קָרַן עוֹר פָּנָיו וַיִּירְאוּ מִגֶּשֶׁת אֵלָיו׃}
}
{
	\eP{
		And Aaron and all the children of Israel saw Moses; and,
		here, the skin of his face was transformed, and they were
		afraid of going over to him.
	}
}
{
}

\Verse{31}
{
	\hP{וַיִּקְרָא אֲלֵהֶם מֹשֶׁה וַיָּשֻׁבוּ אֵלָיו אַהֲרֹן וְכל הַנְּשִׂאִים בָּעֵדָה וַיְדַבֵּר מֹשֶׁה אֲלֵהֶם׃}
}
{
	\eP{
		And Moses called to them. And Aaron and all the chiefs in
		the congregation came back to him, and he spoke to them.
	}
}
{
}

\Verse{32}
{
	\hP{וְאַחֲרֵי כֵן נִגְּשׁוּ כּל בְּנֵי יִשְׂרָאֵל וַיְצַוֵּם אֵת כּל אֲשֶׁר דִּבֶּר יְהֹוָה אִתּוֹ בְּהַר סִינָי׃}
}
{
	\eP{
		And after that all the children of Israel went over. And
		he commanded them everything that YHWH had spoken with him
		in Mount Sinai.
	}
}
{
%	\footnote{
%		after that. First, everyone is afraid when they see Moses. Then, Aaron
%		and the leaders approach Moses. And, after that, the entire people
%		approach him. This indicates the role and importance of leaders. And it
%		indicates that it is natural to be fearful, and that some people
%		(leaders) are needed to be bold, to give others assurance.
%	}
}

\Verse{33}
{
	\hP{וַיְכַל מֹשֶׁה מִדַּבֵּר אִתָּם וַיִּתֵּן עַל פָּנָיו מַסְוֶה׃}
}
{
	\eP{
		And Moses finished speaking with them, and he put a veil
		on his face.
	}
}
{
}

\Verse{34}
{
	\hP{וּבְבֹא מֹשֶׁה לִפְנֵי יְהֹוָה לְדַבֵּר אִתּוֹ יָסִיר אֶת הַמַּסְוֶה עַד צֵאתוֹ וְיָצָא וְדִבֶּר אֶל בְּנֵי יִשְׂרָאֵל אֵת אֲשֶׁר יְצֻוֶּה׃}
}
{
	\eP{
		And when Moses would come in front of YHWH to speak with
		Him, he would turn away the veil until he would go out;
		and he would go out and speak to the children of Israel
		what he had been commanded.
	}
}
{
}

\Verse{35}
{
	\hP{וְרָאוּ בְנֵי יִשְׂרָאֵל אֶת פְּנֵי מֹשֶׁה כִּי קָרַן עוֹר פְּנֵי מֹשֶׁה וְהֵשִׁיב מֹשֶׁה אֶת הַמַּסְוֶה עַל פָּנָיו עַד בֹּאוֹ לְדַבֵּר אִתּוֹ׃}
}
{
	\eP{
		And the children of Israel would see Moses’ face, that the
		skin of Moses’ face was transformed, and Moses would put
		back the veil on his face until he would come to speak
		with Him.
	}
}
{
}
